% ============================================================
%  Annexes — Titre gras centré 24pt Times New Roman
% ============================================================
\clearpage
\thispagestyle{plain}

\vspace*{2cm}

\begin{center}
{\fontsize{24}{29}\selectfont\bfseries Annexes}
\end{center}
\addcontentsline{toc}{chapter}{Annexes}

\vspace{1.5cm}

\section*{A.1~~Structure du projet}

L'arborescence complète du projet \textbf{Dental Reservation} est organisée comme suit~:

\begin{lstlisting}[language={}, caption={Arborescence du projet}]
Dental-reservation/
  api/                          # Backend ASP.NET Core 8
    Controllers/                # 7 controleurs REST
    Models/                     # 8 entites + DTOs
    Data/                       # ApplicationDbContext
    Program.cs                  # Point d'entree
    appsettings.json            # Configuration
    Dockerfile                  # Build multi-stage
  frontend/                     # Frontend React 18
    src/
      pages/                    # 8 pages
      components/               # Composants UI
      hooks/                    # 3 hooks personnalises
      lib/                      # ApiService + utils
    vite.config.ts              # Configuration Vite
    tailwind.config.ts          # Configuration Tailwind
    Dockerfile                  # Build multi-stage + Nginx
  docker-compose.yml            # Orchestration
  diagrams/                     # 8 diagrammes PlantUML
\end{lstlisting}

\section*{A.2~~Configuration Docker Compose}

Le fichier \texttt{docker-compose.yml} orchestre les deux services de l'application~:

\begin{lstlisting}[language=yaml, caption={docker-compose.yml}]
version: "3.8"
services:
  api:
    build:
      context: ./api
      dockerfile: Dockerfile
    ports:
      - "5057:8080"
    environment:
      - ASPNETCORE_ENVIRONMENT=Production
  frontend:
    build:
      context: ./frontend
      dockerfile: Dockerfile
    ports:
      - "80:80"
    depends_on:
      - api
\end{lstlisting}

\section*{A.3~~Fichiers PlantUML}

Les 8 diagrammes UML du projet sont disponibles dans le dossier \texttt{diagrams/}~:

\begin{enumerate}
    \item \texttt{use\_case\_diagram.puml} --- Diagramme de cas d'utilisation
    \item \texttt{class\_diagram.puml} --- Diagramme de classes
    \item \texttt{sequence\_authentification.puml} --- Séquence d'authentification
    \item \texttt{sequence\_rendez\_vous.puml} --- Séquence de gestion des rendez-vous
    \item \texttt{sequence\_patients.puml} --- Séquence de gestion des patients
    \item \texttt{activity\_authentification.puml} --- Activité d'authentification
    \item \texttt{activity\_rendez\_vous.puml} --- Activité de création d'un rendez-vous
    \item \texttt{activity\_patients.puml} --- Activité de gestion des patients (CRUD)
\end{enumerate}

\clearpage
